\section*{Capítulo \ref{ch:b1:field-lab-methods-1} --- Métodos de campo e de laboratório I}

\begin{solution}{exo:b1:field-lab-methods-1:1}
Na escala 1:10\,000, \qty{4.9}{m} tornam-se $4.9/10\,000 = \qty{0.49}{mm}$; na
escala 1:25\,000, $4.9/25\,000 = \qty{0.20}{mm}$, cerca da largura de um traço de lápis.
\end{solution}

\begin{solution}{exo:b1:field-lab-methods-1:2}
Em Reykjavik $D = \ang{-11.28}$: $52 - 11.28 = \ang{040.7}$. Em Wellington
$D = \ang{+23.23}$: $52 + 23.23 = \ang{075.2}$. A mesma leitura de bússola significa
direções separadas por \ang{34.5}.
\end{solution}

\begin{solution}{exo:b1:field-lab-methods-1:3}
$20\to015$. Com $t = \ang{15}$ e $p = \ang{20}$, o vetor unitário para baixo é
\[
  (\cos\ang{20}\sin\ang{15},\ \cos\ang{20}\cos\ang{15},\ -\sin\ang{20}) =
  (0.2432,\ 0.9077,\ -0.3420).
\]
\end{solution}

\begin{solution}{exo:b1:field-lab-methods-1:4}
A média aritmética é $(350 + 358 + 6)/3 = \ang{238}$, quase oposta às três
leituras. Os vetores unitários horizontais $(\sin a, \cos a)$ somam-se em
$(-0.1736 - 0.0349 + 0.1045,\ 0.9848 + 0.9994 + 0.9945) = (-0.1040, 2.9787)$,
cujo \omterm{def:b1:spherical-astronomy:horizontal}{azimute} é $\operatorname{atan2}(-0.1040, 2.9787) = \ang{-2.0}$, isto é,
\ang{358.0}: a média correta, no meio das leituras.
\end{solution}

\begin{solution}{exo:b1:field-lab-methods-1:5}
Média $175/5 = \ang{35.0}$. Desvios $-1, 2, 0, -2, 1$; soma dos quadrados 10;
$s = \sqrt{10/4} = \ang{1.58}$. Erro padrão $1.58/\sqrt5 = \ang{0.71}$.
\end{solution}

\begin{solution}{exo:b1:field-lab-methods-1:6}
As normais são $(\sin\ang{30}, 0, \cos\ang{30}) = (0.5, 0, 0.866)$ e
$(-0.5, 0, 0.866)$; $\vec S = (0, 0, 1.732)$ é vertical: a média é um plano horizontal, e $\bar R = 1.732/2 = 0.866$, cada leitura a \ang{30} da média. As duas leituras não amostram uma superfície com erros: pertencem a dois flancos de uma \omterm{def:b1:geological-maps:fold}{dobra}, duas populações diferentes, e a sua média não descreve nenhuma delas. Cada flanco é calculado em média separadamente; juntos eles definem o eixo da \omterm{def:b1:geological-maps:fold}{dobra}, a reta comum aos dois planos, aqui horizontal e norte--sul.
\end{solution}

\begin{solution}{exo:b1:field-lab-methods-1:7}
$\tan p = \tan\ang{12}\cos(\ang{80} - \ang{120}) = 0.2126\times0.7660 =
0.1628$: $p = \ang{9.2}$, com as estrias caindo em direção a \ang{120}. Para o
\omterm{def:b1:field-lab-methods-1:trend-plunge}{rumo} \ang{170}, $\cos(\ang{-90}) = 0$: a linha é horizontal, ao longo da \omterm{def:b1:geological-maps:attitude}{direção}.
\end{solution}

\begin{solution}{exo:b1:field-lab-methods-1:8}
$0.47/0.20 = 2.35$ anos depois de 1º de janeiro de 2026, isto é, na primavera de 2028. Em 1º de janeiro de 2030: $-0.47 + 4\times0.20 = \ang{+0.33}$, a leste do norte geográfico.
\end{solution}

\begin{solution}{exo:b1:field-lab-methods-1:9}
As três normais estão no plano vertical de \omterm{def:b1:spherical-astronomy:horizontal}{azimute} \ang{080}, a \ang{10}, \ang{12} e \ang{14} da vertical; as duas externas são simétricas em torno da do meio, de modo que a sua soma fica sobre ela: a média é
$080/12$, com $\theta = \ang{2}, 0, \ang{2}$. Então $\bar R = (1 +
2\cos\ang{2})/3 = (1 + 1.998782)/3 = 0.999594$, e $\sqrt{2(1 - \bar R)} =
\sqrt{0.000812} = \qty{0.02850}{rad} = \ang{1.633}$; exatamente,
$\theta_{\text{rms}} = \sqrt{(4 + 0 + 4)/3} = \ang{1.633}$.
\end{solution}

\begin{solution}{exo:b1:field-lab-methods-1:10}
A seção é um círculo de raio $\sqrt{r^2 - h^2}$. Sua média é
$\frac1r\int_0^r\sqrt{r^2 - h^2}\,\mathrm dh = \frac1r\cdot\frac{\pi r^2}{4}
= \frac{\pi r}{4}$, pois a integral é a área de um quarto de disco. Grãos
de \qty{0.50}{mm} de diâmetro aparecem, em média, com $0.50\times\pi/4 = \qty{0.39}{mm}$
de diâmetro: a \omterm{def:b1:field-lab-methods-1:thin-section}{lâmina delgada} subestima os tamanhos de grão.
\end{solution}

\begin{solution}{exo:b1:field-lab-methods-1:11}
A rotação de $D$ em torno da vertical,
$(x, y, z) \mapsto (x\cos D + y\sin D,\ -x\sin D + y\cos D,\ z)$, leva
$(\sin\delta\sin\alpha, \sin\delta\cos\alpha, \cos\delta)$ a
$(\sin\delta\sin(\alpha + D), \sin\delta\cos(\alpha + D), \cos\delta)$: ela
soma $D$ a cada \omterm{def:b1:geological-maps:attitude}{sentido de mergulho}. Sendo linear, leva a resultante
$\vec S$ à resultante das normais giradas; ela conserva os comprimentos e a componente vertical, de modo que o \omterm{def:b1:geological-maps:attitude}{mergulho} médio $\arccos(S_U/\|\vec S\|)$ não muda, e gira a parte horizontal de $\vec S$ de $D$, de modo que o \omterm{def:b1:geological-maps:attitude}{sentido de mergulho} médio aumenta de $D$. Corrigir cada leitura pela declinação, ou apenas a sua média, dá o mesmo resultado.
\end{solution}

\begin{solution}{exo:b1:field-lab-methods-1:12}
A dispersão ao longo da linha de maior \omterm{def:b1:geological-maps:attitude}{mergulho} é uma dispersão do \omterm{def:b1:geological-maps:attitude}{mergulho}: o erro padrão é $2/\sqrt N \le 0.5$ para $N \ge 16$. Transversalmente, uma rotação da normal de \ang{2} muda o \omterm{def:b1:geological-maps:attitude}{sentido de mergulho} de $2/\sin\ang{12} = \ang{9.62}$;
o erro padrão do \omterm{def:b1:geological-maps:attitude}{sentido de mergulho} médio é $9.62/\sqrt N \le 1$ para
$N \ge 92.5$, isto é, 93 leituras. Um \omterm{def:b1:geological-maps:attitude}{mergulho} suave se conhece depressa; a sua \omterm{def:b1:geological-maps:attitude}{direção}, não.
\end{solution}

\begin{solution}{pb:b1:field-lab-methods-1:1}
\textbf{1.} \omterm{def:b1:geological-maps:attitude}{Sentidos de mergulho}: $967/12 = \ang{80.6}$; \omterm{def:b1:geological-maps:attitude}{mergulhos}: $145/12 =
\ang{12.1}$.

\textbf{2.} \ang{355} e \ang{005} têm média \ang{180}, a direção oposta. Aqui todos os \omterm{def:b1:geological-maps:attitude}{sentidos de mergulho} estão entre \ang{076} e \ang{085}, longe do salto em \ang{000}, de modo que a média aritmética funciona.

\textbf{3.} $\|\vec n\|^2 = \sin^2\delta(\sin^2\alpha + \cos^2\alpha) +
\cos^2\delta = 1$. A linha de \omterm{def:b1:geological-maps:attitude}{mergulho} para baixo tem a \omterm{def:b1:geological-maps:attitude}{direção}
$(\cos\delta\sin\alpha, \cos\delta\cos\alpha, -\sin\delta)$, e o seu produto
escalar com $\vec n$ é $\sin\delta\cos\delta - \cos\delta\sin\delta = 0$; a
linha de \omterm{def:b1:geological-maps:attitude}{direção} é horizontal e perpendicular ao \omterm{def:b1:geological-maps:attitude}{sentido de mergulho},
$(\cos\alpha, -\sin\alpha, 0)$, e o seu produto escalar com $\vec n$ é
$\sin\delta\sin\alpha\cos\alpha - \sin\delta\cos\alpha\sin\alpha = 0$.

\textbf{4.} A leitura 11 é a mesma que a 2: $(0.2222, 0.0352, 0.9744)$.
A leitura 12:
\[
  (\sin\ang{11}\sin\ang{82},\ \sin\ang{11}\cos\ang{82},\ \cos\ang{11}) =
  (0.1890,\ 0.0266,\ 0.9816).
\]

\textbf{5.} $\vec S = (2.4747,\ 0.4082,\ 11.7294)$.

\textbf{6.} $\|\vec S\| = \sqrt{6.1241 + 0.1666 + 137.5788} = 11.99457$, e
$\bar R = 11.99457/12 = 0.999547$.

\textbf{7.} $\bar\alpha = \operatorname{atan2}(2.4747, 0.4082) = \ang{80.6}$
e $\bar\delta = \arccos(11.7294/11.99457) = \arccos 0.977893 = \ang{12.1}$:
a média vetorial é $080.6/12.1$, a mesma das médias aritméticas aqui.

\textbf{8.} \omterm{def:b1:geological-maps:attitude}{Sentidos de mergulho}: $s = \ang{2.94}$; \omterm{def:b1:geological-maps:attitude}{mergulhos}: $s = \ang{1.68}$.

\textbf{9.} $\theta_{\text{rms}} \approx \sqrt{2\times0.000453} =
\qty{0.0301}{rad} = \ang{1.72}$.

\textbf{10.} $\sqrt{1.68^2 + (\sin\ang{12.1}\times2.94)^2} = \sqrt{2.82 +
0.38} = \ang{1.79}$, próximo de \ang{1.72}: as leituras se comportam como o modelo de erro prevê, dentro dos efeitos do arredondamento ao grau e de uma pequena amostra.

\textbf{11.} O \omterm{def:b1:geological-maps:attitude}{mergulho}: ele contribui com 2.82 dos 3.20 (em graus quadrados).
O \omterm{def:b1:geological-maps:attitude}{sentido de mergulho} conta apenas por meio de $\sin\delta = 0.21$, porque num estrato suave uma mudança do \omterm{def:b1:geological-maps:attitude}{sentido de mergulho} gira o plano muito pouco.

\textbf{12.} \omterm{def:b1:geological-maps:attitude}{Mergulho} médio: $1.68/\sqrt{12} = \ang{0.48}$; \omterm{def:b1:geological-maps:attitude}{sentido de mergulho} médio:
$2.94/\sqrt{12} = \ang{0.85}$.

\textbf{13.} Não: \ang{2.9} de \omterm{def:b1:geological-maps:attitude}{sentido de mergulho} num estrato que mergulha \ang{12} são apenas
$2.9\sin\ang{12.1} = \ang{0.6}$ de giro do plano, menos que os \ang{1.7} dos \omterm{def:b1:geological-maps:attitude}{mergulhos}. Qualquer leve ondulação de uma superfície de acamamento suave espalha amplamente o seu \omterm{def:b1:geological-maps:attitude}{sentido de mergulho}.

\textbf{14.} $D = \ang{-2.5}$: $83.1 - 2.5 = \ang{080.6}$; a sua média é
$080.6/12.1$, a mesma do primeiro geólogo.

\textbf{15.} O \omterm{def:b1:geological-maps:attitude}{mergulho} é lido com o clinômetro, em relação à vertical dada pela gravidade; a agulha não tem papel nisso.

\textbf{16.} Sim: uma rotação em torno da vertical desloca toda leitura e a média do mesmo ângulo (\cref{exo:b1:field-lab-methods-1:11}); a resposta é a mesma.

\textbf{17.} $10\times0.2 = \ang{2}$ de erro em cada \omterm{def:b1:geological-maps:attitude}{sentido de mergulho}, mais que o dobro do erro padrão de \ang{0.85}: um erro sistemático que a média não consegue detectar.

\textbf{18.} $80.6 - 90 = \ang{350.6}$; a \omterm{def:b1:geological-maps:attitude}{atitude} é \omterm{def:b1:geological-maps:attitude}{direção} \ang{351}, \omterm{def:b1:geological-maps:attitude}{mergulho} \ang{12} para leste.

\textbf{19.} \omterm{def:b1:geological-maps:attitude}{Sentidos de mergulho} $080.6$ e $080.6$: nenhuma diferença com esta precisão; \omterm{def:b1:geological-maps:attitude}{mergulhos} $12.1$ e $12.2$: \ang{0.1}.

\textbf{20.} As normais $(\sin\ang{12.2}\sin\ang{80.6},
\sin\ang{12.2}\cos\ang{80.6}, \cos\ang{12.2})$ e $\vec S/\|\vec S\|$ fazem
um ângulo de \ang{0.13}.

\textbf{21.} $1.72/\sqrt{12} = \ang{0.50}$; a diferença, \ang{0.13}, é um quarto disso: não significativa. A bússola e as sondagens concordam.

\textbf{22.} Quatro vezes mais: 48 leituras.

\textbf{23.} Por exemplo: um ajuste de declinação errado ou desatualizado; uma bússola desviada por ferro próximo (um veículo, ferramentas, rochas magnéticas); um clinômetro desajustado; leituras feitas só nas superfícies mais expostas, que podem não representar o flanco. A média reduz os erros aleatórios, e não os sistemáticos.

\textbf{24.} A \omterm{def:b1:geological-maps:attitude}{atitude} média do flanco leste é \textbf{081/12}: \omterm{def:b1:geological-maps:attitude}{sentido de mergulho}
\ang{081}, \omterm{def:b1:geological-maps:attitude}{mergulho} \ang{12}.
\end{solution}
